\section{Introduction}
\label{sec:intro}
Mixture-of-experts (MoE) models route each token to selected experts distributed across GPUs. Dispatch sends activations to expert ranks; Combine returns their outputs to token owners. DeepEP and NCCL EP optimize these transfers using GPU-initiated transport, staging, and topology-aware data movement~\cite{deepep,ncclep}. At small batches, however, shortening the payload phase leaves the placement of control operations increasingly visible. A sender can have data ready while waiting for peer readiness, and a receiver can wait for an endpoint-generated completion signal after the sender's transfer work has already progressed.

Our primary scenario is \emph{cross-node DeepEP V2 Direct Dispatch and Combine}. Both contain an entry barrier and a separate tail barrier~\cite{deepepcode}. The entry barrier precedes payload injection; the tail coordinates sender completion and local workers before exchanging completion signals across the EP group. These boundaries can delay both the start of data transfer and the release of its consumer.

We present \sys{}, which places both coordination decisions on the managed network path. At the head, each rank announces readiness and immediately uploads \data{} into reserved network staging. The \sys node buffers early arrivals until the synchronization domain is ready, then streams buffered and arriving data to the destinations. At the tail, the node combines a sealed traffic plan for each destination with output progress and appends an ordered \done{} after that destination's last data. This establishes destination input completion without waiting for a sender to return through its software completion and local coordination path to initiate a new signal exchange.

The two mechanisms act on different dependencies. The first overlaps peer-readiness aggregation with upload. The second changes where completion notification is generated. The receiver must still obtain its complete, visible input; the sender must still protect buffers that the transport may read. Generation-specific ownership and separate recycling keep these requirements distinct. Existing multicast and reduction can remain in the data plane, but their arithmetic and byte savings are not part of the synchronization claim.

Following Swift~\cite{swift}, we first decompose delays into endpoint and fabric components, then compare their causal placement. A symmetric reference exposing a separate one-way traversal on each side of \data{} can save approximately half a fabric RTT at each boundary. For actual Direct execution, the relevant tail quantity is the interval between data visibility and release through the tail signals. The duration of this interval depends on transfer progress and notification timing. Our model retains sender completion time, output service, and notification overhead so that this distinction is explicit.

Measurements on two H200 nodes motivate the small-batch focus. On cross-node Direct at eight tokens per rank, the entry exchange occupies about 20\% of the traced Dispatch operator and the completion-signal window a further 19\%; raising the batch sixteen-fold leaves both windows at 12--15\us{} while their shares fall below 5\%. \secref{sec:smallbatch} explains why the denominator, not the coordination, is what changes, and \secref{sec:evidence} reports the measurements in full.

We contribute an analysis of Direct's two coordination boundaries, two in-network mechanisms that overlap readiness and generate completion at the network output, and a common delay model with explicit visibility and resource conditions. Throughout, the baselines are measured and the latency opportunity is analytical: we build no in-network prototype and claim no speedup.
